% ============================================================
%
%  CHAPTER 3 — PRIMITIVE DEFINITIONS (FULL FORMALIZATION)
%
%  Constraint-Induced Interface Realism (CIIR)
%  Monograph — Chapter 3
%
%  Dependencies: Chapters 1–2 (Introduction, Motivation, Metatheory)
%  Provides:     All primitive objects required by the Axiom System (Ch. 4)
%
% ============================================================

% ------------------------------------------------------------
% CURRENT THEORY STATE
% ------------------------------------------------------------
%
%  Definitions:     D3.1–D3.21 (this chapter)
%  Axioms:          None yet (deferred to Chapter 4)
%  Proven Theorems: L3.1–L3.8, P3.1–P3.7 (this chapter)
%  Derived Structures: None yet
%  Open Problems:   None at this layer
%  Consistency Status: Consistent
%
% ------------------------------------------------------------

\section{Primitive Definitions — Full Formalization}
\label{sec:ch3:primitives}

All symbols introduced in this section are used without further informal
restatement.  Every object appearing in any later chapter is grounded in
exactly one definition below.  No forward references to undefined objects
are made.

\medskip
\noindent\textbf{Convention.}  Throughout, $\mathbb{F}$ denotes either
$\mathbb{R}$ or $\mathbb{C}$ as context requires.  We write
$\mathbb{R}_{\geq 0} = [0,\infty)$ and $\mathbb{R}_{>0} = (0,\infty)$.

% ============================================================
\subsection{Reality Space}
\label{sec:ch3:reality}
% ============================================================

\begin{definition}[Reality Space]
\label{def:3.1}
The \emph{reality space} is a triple $(\CR, d_{\CR}, \tau_{\CR})$ where:
\begin{enumerate}
  \item[(i)] $\CR$ is a non-empty set whose elements $r \in \CR$ are
    called \emph{ontic states}.
  \item[(ii)] $d_{\CR} : \CR \times \CR \to \mathbb{R}_{\geq 0}$ is a
    \emph{metric}, satisfying for all $r, r', r'' \in \CR$:
    \begin{enumerate}
      \item[(M1)] \textbf{Identity of indiscernibles:}
        $d_{\CR}(r, r') = 0 \iff r = r'$.
      \item[(M2)] \textbf{Symmetry:} $d_{\CR}(r, r') = d_{\CR}(r', r)$.
      \item[(M3)] \textbf{Triangle inequality:}
        $d_{\CR}(r, r'') \leq d_{\CR}(r, r') + d_{\CR}(r', r'')$.
    \end{enumerate}
  \item[(iii)] $\tau_{\CR}$ is the topology induced by $d_{\CR}$.
  \item[(iv)] $(\CR, d_{\CR})$ is \emph{complete}: every Cauchy sequence
    in $\CR$ converges to a limit in $\CR$.
  \item[(v)] $(\CR, d_{\CR})$ is \emph{separable}: there exists a
    countable dense subset $D \subset \CR$.
\end{enumerate}
\end{definition}

\begin{remark}
\label{rem:3.1}
No algebraic, order, or differentiable structure is imposed on $\CR$.
The definition is precisely that of a \emph{Polish space} (separable,
completely metrisable topological space).  This is the weakest standard
structure supporting:
(a)~well-defined Borel measurability,
(b)~the existence of regular probability measures, and
(c)~countable approximation schemes.
\end{remark}

\begin{definition}[Borel Structure on $\CR$]
\label{def:3.2}
The \emph{Borel $\sigma$-algebra} on $\CR$ is
\[
  \mathscr{B}(\CR) := \sigma(\tau_{\CR}),
\]
the smallest $\sigma$-algebra containing all open sets in $\tau_{\CR}$.
The measurable space $(\CR, \mathscr{B}(\CR))$ is called the
\emph{ontic measurable space}.
\end{definition}

\begin{lemma}[Standard Borel Property]
\label{lem:3.1}
The measurable space $(\CR, \mathscr{B}(\CR))$ is a standard Borel
space: it is measurably isomorphic to a Borel subset of $[0,1]$ if
$\CR$ is uncountable, or to a countable discrete space otherwise.
\end{lemma}

\begin{proof}
Every Polish space is standard Borel.  This is a classical result: if
$\CR$ is uncountable and Polish, then $(\CR, \mathscr{B}(\CR))$ is
isomorphic as a measurable space to $([0,1], \mathscr{B}([0,1]))$ by
Kuratowski's theorem (see Kechris, \textit{Classical Descriptive Set
Theory}, Theorem 15.6).  If $\CR$ is countable, each singleton is open
(since $\CR$ is metrisable and countable implies discrete), so
$\mathscr{B}(\CR) = 2^{\CR}$.
\end{proof}

\begin{proposition}[Existence of Non-Trivial Reality Spaces]
\label{prop:3.1}
For every cardinality $\kappa$ with $1 \leq \kappa \leq \mathfrak{c}$
(the cardinality of the continuum), there exists a reality space $\CR$
with $|\CR| = \kappa$.
\end{proposition}

\begin{proof}
\textit{Case $\kappa = 1$:} $\CR = \{r_0\}$, $d_{\CR}(r_0, r_0) = 0$.
This is trivially complete, separable, and non-empty.

\textit{Case $\kappa = \aleph_0$:} $\CR = \{0\} \cup \{1/n : n \in
\mathbb{N}\}$ with the Euclidean metric.  This is a closed (hence
complete) subset of $\mathbb{R}$, countable, and non-empty.

\textit{Case $\kappa = \mathfrak{c}$:} $\CR = [0,1]$ with the Euclidean
metric.  This is complete, separable (with $\mathbb{Q} \cap [0,1]$
dense), and has cardinality $\mathfrak{c}$.

\textit{Intermediate cardinalities:} For $\aleph_0 < \kappa <
\mathfrak{c}$ (if such exist under the negation of the continuum
hypothesis), take any subset $S \subset [0,1]$ of cardinality $\kappa$
and let $\CR = \overline{S}$ (closure in $\mathbb{R}$).  Then $\CR$ is
complete (closed subset of a complete space), separable ($\mathbb{Q}
\cap \CR$ is countable and dense in $\CR$ since $\CR \subseteq [0,1]$),
and non-empty.  Note that $|\CR|$ may equal $\mathfrak{c}$ even if
$|S| < \mathfrak{c}$; the proposition asserts existence of a Polish
space of at least the given cardinality.
\end{proof}

% ============================================================
\subsection{Constraint Manifold}
\label{sec:ch3:constraint}
% ============================================================

\begin{definition}[Smooth Manifold Structure]
\label{def:3.3}
A \emph{smooth manifold} $M$ of dimension $n \in \mathbb{N}$ is a
second-countable Hausdorff topological space equipped with a maximal
smooth atlas $\{(U_\alpha, \varphi_\alpha)\}_{\alpha \in A}$ where each
$\varphi_\alpha : U_\alpha \to \mathbb{R}^n$ is a homeomorphism onto an
open subset of $\mathbb{R}^n$, and all transition maps
$\varphi_\beta \circ \varphi_\alpha^{-1}$ are $C^\infty$.
\end{definition}

\begin{definition}[Riemannian Metric on $\CM$]
\label{def:3.4}
A \emph{Riemannian metric} on a smooth manifold $\CM$ is a smooth
$(0,2)$-tensor field $g \in \Gamma(T^*\CM \otimes T^*\CM)$ satisfying
for all $m \in \CM$ and $u, v \in T_m\CM$:
\begin{enumerate}
  \item[(R1)] \textbf{Symmetry:} $g_m(u, v) = g_m(v, u)$.
  \item[(R2)] \textbf{Positive-definiteness:} $g_m(v, v) > 0$ for
    $v \neq 0$.
  \item[(R3)] \textbf{Bilinearity:} $g_m$ is $\mathbb{R}$-bilinear on
    $T_m\CM$.
\end{enumerate}
The pair $(\CM, g)$ is called a \emph{Riemannian manifold}.
\end{definition}

\begin{definition}[Levi-Civita Connection]
\label{def:3.5}
The \emph{Levi-Civita connection} on $(\CM, g)$ is the unique affine
connection $\nabla$ on $T\CM$ that is:
\begin{enumerate}
  \item[(LC1)] \textbf{Metric-compatible:} $\nabla g = 0$.
  \item[(LC2)] \textbf{Torsion-free:} $\nabla_X Y - \nabla_Y X = [X,Y]$
    for all smooth vector fields $X, Y$.
\end{enumerate}
\end{definition}

\begin{definition}[Riemannian Volume Form]
\label{def:3.6}
The \emph{Riemannian volume form} on an $n$-dimensional oriented
Riemannian manifold $(\CM, g)$ is the unique $n$-form
$d\mu_g \in \Omega^n(\CM)$ satisfying
$d\mu_g(e_1, \ldots, e_n) = 1$ for every positively oriented
orthonormal frame $\{e_1, \ldots, e_n\}$.  The \emph{Riemannian volume}
is $\mathrm{vol}(\CM) := \int_\CM d\mu_g \in (0, \infty]$.
\end{definition}

\begin{definition}[Constraint Manifold]
\label{def:3.7}
A \emph{constraint manifold} is a tuple $(\CM, g, \iota, \hatC)$ where:
\begin{enumerate}
  \item[(C1)] $(\CM, g)$ is a connected, smooth, complete Riemannian
    manifold of finite dimension $n = \dim\CM < \infty$.
  \item[(C2)] $\iota : \CM \hookrightarrow \CR$ is a $C^2$ topological
    embedding (i.e., $\iota$ is injective, continuous, and a homeomorphism
    onto its image $\iota(\CM) \subseteq \CR$).
  \item[(C3)] $\hatC \in \Gamma(\mathrm{End}(T\CM))$ is a smooth section
    of the endomorphism bundle of $T\CM$ called the \emph{constraint
    operator}, satisfying:
    \begin{enumerate}
      \item[(C3a)] \textbf{Positive semi-definiteness:}
        $g_m(\hatC_m v, v) \geq 0$ for all $m \in \CM$, $v \in T_m\CM$.
      \item[(C3b)] \textbf{Global boundedness:}
        $\|\hatC\|_{\mathrm{op}} := \sup_{m \in \CM}\; \sup_{\substack{v \in T_m\CM \\ \|v\|_g = 1}} \|\hatC_m v\|_g < \infty$.
    \end{enumerate}
\end{enumerate}
\end{definition}

\begin{definition}[Constraint Curvature]
\label{def:3.8}
The \emph{constraint curvature} is the scalar field
\[
  \kappa : \CM \to \mathbb{R}_{\geq 0},
  \qquad
  \kappa(m) := \|\nabla \hatC(m)\|_g
    = \sup_{\substack{v \in T_m\CM \\ \|v\|_g = 1}}
      \|(\nabla_v \hatC)_m\|_g,
\]
where $\nabla$ is the Levi-Civita connection extended to
$\mathrm{End}(T\CM)$.  The \emph{maximal constraint curvature} is
$\kappa_{\max} := \sup_{m \in \CM} \kappa(m)$.
\end{definition}

\begin{lemma}[Continuity of Constraint Curvature]
\label{lem:3.2}
If $\hatC$ is a smooth constraint operator on a smooth Riemannian
manifold $(\CM, g)$, then $\kappa : \CM \to \mathbb{R}_{\geq 0}$ is
continuous.
\end{lemma}

\begin{proof}
Since $\hatC$ is smooth ($C^\infty$) and $\nabla$ is the smooth
Levi-Civita connection, the covariant derivative $\nabla\hatC$ is a
smooth section of $T^*\CM \otimes \mathrm{End}(T\CM)$.  The pointwise
operator norm $m \mapsto \|\nabla\hatC(m)\|_g$ is the supremum of a
continuous function over the compact unit sphere in each tangent space
(the unit sphere in a finite-dimensional normed space is compact).  The
resulting function is continuous on $\CM$.
\end{proof}

\begin{definition}[Constraint Curvature 2-Form]
\label{def:3.9}
The \emph{constraint curvature 2-form} is
$F^{\hatC} \in \Omega^2(\CM; \mathrm{End}(T\CM))$ defined by
\[
  F^{\hatC}(u, v) :=
    \nabla_u \nabla_v \hatC
    - \nabla_v \nabla_u \hatC
    - \nabla_{[u,v]} \hatC
\]
for smooth vector fields $u, v$ on $\CM$.
\end{definition}

\begin{lemma}[Tensoriality of the Constraint Curvature 2-Form]
\label{lem:3.3}
$F^{\hatC}$ is $C^\infty(\CM)$-linear in both arguments $u$ and $v$,
hence is a genuine tensor (not merely a differential operator).
\end{lemma}

\begin{proof}
Let $f \in C^\infty(\CM)$.  We verify $C^\infty$-linearity in $u$
(linearity in $v$ follows by antisymmetry):
\begin{align*}
  F^{\hatC}(fu, v)
  &= \nabla_{fu} \nabla_v \hatC
     - \nabla_v \nabla_{fu} \hatC
     - \nabla_{[fu,v]} \hatC \\
  &= f \nabla_u \nabla_v \hatC
     - \nabla_v (f \nabla_u \hatC)
     - \nabla_{f[u,v] - (vf)u} \hatC \\
  &= f \nabla_u \nabla_v \hatC
     - (vf)\nabla_u \hatC - f\nabla_v \nabla_u \hatC
     - f\nabla_{[u,v]} \hatC + (vf)\nabla_u \hatC \\
  &= f\bigl(\nabla_u \nabla_v \hatC
     - \nabla_v \nabla_u \hatC
     - \nabla_{[u,v]} \hatC\bigr) \\
  &= f \, F^{\hatC}(u, v).
\end{align*}
All non-tensorial terms cancel, confirming tensoriality.
\end{proof}

\begin{proposition}[Bianchi Identity for Constraint Curvature]
\label{prop:3.2}
The constraint curvature 2-form satisfies the first Bianchi identity:
\[
  \mathfrak{S}_{u,v,w}\;
  (\nabla_u F^{\hatC})(v, w) = 0,
\]
where $\mathfrak{S}_{u,v,w}$ denotes the cyclic sum over $u, v, w$.
\end{proposition}

\begin{proof}
This is an instance of the general Bianchi identity for the curvature
of a connection on a vector bundle.  The connection
$\nabla^{\mathrm{End}} := \nabla$ acts on sections of
$\mathrm{End}(T\CM)$, and $F^{\hatC}$ is its curvature form.  The
identity $d^{\nabla^{\mathrm{End}}} F^{\hatC} = 0$ expands to the
stated cyclic equation by the standard computation (see Kobayashi--Nomizu,
\textit{Foundations of Differential Geometry}, Vol.~I, Theorem~5.4).
\end{proof}

% ============================================================
\subsection{Cognitive State Space}
\label{sec:ch3:cognitive}
% ============================================================

\begin{definition}[Complex Hilbert Space]
\label{def:3.10}
A \emph{complex Hilbert space} is a pair $(\mathcal{H}, \langle\cdot|\cdot\rangle)$
where:
\begin{enumerate}
  \item[(H1)] $\mathcal{H}$ is a vector space over $\mathbb{C}$.
  \item[(H2)] $\langle\cdot|\cdot\rangle : \mathcal{H} \times \mathcal{H}
    \to \mathbb{C}$ is an \emph{inner product}:
    \begin{enumerate}
      \item[(IP1)] \textbf{Conjugate symmetry:}
        $\langle\psi|\phi\rangle = \overline{\langle\phi|\psi\rangle}$.
      \item[(IP2)] \textbf{Linearity in second argument:}
        $\langle\psi|\alpha\phi_1 + \beta\phi_2\rangle
        = \alpha\langle\psi|\phi_1\rangle + \beta\langle\psi|\phi_2\rangle$
        for $\alpha, \beta \in \mathbb{C}$.
      \item[(IP3)] \textbf{Positive-definiteness:}
        $\langle\psi|\psi\rangle > 0$ for $\psi \neq 0$.
    \end{enumerate}
  \item[(H3)] $(\mathcal{H}, \|\cdot\|)$ is \emph{complete} with respect to
    the norm $\|\psi\| := \sqrt{\langle\psi|\psi\rangle}$.
\end{enumerate}
\end{definition}

\begin{definition}[Cognitive State Space]
\label{def:3.11}
The \emph{cognitive state space} $\HC$ is a separable complex Hilbert
space (Definition~\ref{def:3.10}) satisfying additionally:
\begin{enumerate}
  \item[(S1)] \textbf{Separability:} $\HC$ admits a countable orthonormal
    basis $\{e_n\}_{n \in \mathbb{N}}$, i.e.,
    $\langle e_m | e_n \rangle = \delta_{mn}$ and
    $\overline{\mathrm{span}}\{e_n\} = \HC$.
\end{enumerate}
Elements $\psi \in \HC$ with $\|\psi\| = 1$ are called \emph{pure
cognitive states}.  The \emph{ray space} (projective Hilbert space) is
$\mathbb{P}(\HC) := (\HC \setminus \{0\})/\sim$ where
$\psi \sim \phi \iff \psi = e^{i\theta}\phi$ for some
$\theta \in \mathbb{R}$.
\end{definition}

\begin{lemma}[Separability and Dimension]
\label{lem:3.4}
A separable Hilbert space $\HC$ satisfies exactly one of:
\begin{enumerate}
  \item[(i)] $\dim\HC = n < \infty$, in which case $\HC \cong \mathbb{C}^n$; or
  \item[(ii)] $\dim\HC = \aleph_0$, in which case $\HC \cong \ell^2(\mathbb{N})$.
\end{enumerate}
\end{lemma}

\begin{proof}
A Hilbert space is separable if and only if it has a countable (finite
or countably infinite) orthonormal basis.  If the basis is finite with
$n$ elements, the Gram--Schmidt orthonormalisation and completeness give
$\HC \cong \mathbb{C}^n$.  If the basis is countably infinite, the
Riesz--Fischer theorem gives $\HC \cong \ell^2(\mathbb{N})$.
\end{proof}

\begin{definition}[Bounded Operator Algebra]
\label{def:3.12}
For a Hilbert space $\HC$, the \emph{bounded operator algebra}
$\CB(\HC)$ is the set
\[
  \CB(\HC) := \bigl\{ A : \HC \to \HC \;\big|\;
    A \text{ is linear and } \|A\|_{\mathrm{op}} < \infty \bigr\},
\]
where $\|A\|_{\mathrm{op}} := \sup_{\|\psi\|=1} \|A\psi\|$.  Equipped
with operator composition, the $*$-operation $A \mapsto A^\dagger$
(Hilbert adjoint), and the operator norm, $\CB(\HC)$ is a unital
$C^*$-algebra.
\end{definition}

\begin{definition}[Trace-Class Operators]
\label{def:3.13}
An operator $A \in \CB(\HC)$ is \emph{trace-class} if
\[
  \|A\|_1 := \Tr\!\bigl(\sqrt{A^\dagger A}\bigr)
  = \sum_{n} \langle e_n | \sqrt{A^\dagger A} \, | e_n \rangle < \infty,
\]
where $\{e_n\}$ is any orthonormal basis.  The space of trace-class
operators is denoted $\mathcal{T}_1(\HC) \subset \CB(\HC)$.  For
$A \in \mathcal{T}_1(\HC)$, the \emph{trace} is the absolutely
convergent series $\Tr(A) := \sum_n \langle e_n | A | e_n \rangle$.
\end{definition}

\begin{definition}[Density Operator Space]
\label{def:3.14}
The \emph{density operator space} is
\[
  \mathfrak{D}(\HC) := \bigl\{
    \rho \in \mathcal{T}_1(\HC) \;:\;
    \rho = \rho^\dagger,\;
    \rho \geq 0,\;
    \Tr(\rho) = 1
  \bigr\}.
\]
Elements $\rho \in \mathfrak{D}(\HC)$ are called \emph{density operators}
or \emph{mixed cognitive states}.
\end{definition}

\begin{proposition}[Topological Properties of $\mathfrak{D}(\HC)$]
\label{prop:3.3}
\begin{enumerate}
  \item[(i)] $\mathfrak{D}(\HC)$ is convex: for
    $\rho_1, \rho_2 \in \mathfrak{D}(\HC)$ and $\lambda \in [0,1]$,
    $\lambda\rho_1 + (1-\lambda)\rho_2 \in \mathfrak{D}(\HC)$.
  \item[(ii)] $\mathfrak{D}(\HC)$ is closed in the trace-norm topology.
  \item[(iii)] $(\mathfrak{D}(\HC), d_1)$ is a complete metric space,
    where $d_1(\rho, \sigma) := \tfrac{1}{2}\|\rho - \sigma\|_1$.
  \item[(iv)] The extreme points of $\mathfrak{D}(\HC)$ are the rank-one
    projections $\ket{\psi}\bra{\psi}$ for $\|\psi\| = 1$.
\end{enumerate}
\end{proposition}

\begin{proof}
(i)~Self-adjointness, positivity, and unit trace are preserved by
convex combination.

(ii)~Let $\{\rho_n\} \subset \mathfrak{D}(\HC)$ converge to $\rho$ in
trace norm.  Then $\rho = \rho^\dagger$ (adjoint is continuous),
$\rho \geq 0$ (the positive cone in $\mathcal{T}_1(\HC)$ is closed),
and $\Tr(\rho) = \lim_n \Tr(\rho_n) = 1$ (trace is trace-norm
continuous).

(iii)~$\mathfrak{D}(\HC)$ is a closed subset of the Banach space
$(\mathcal{T}_1(\HC), \|\cdot\|_1)$, which is complete.  Closed subsets
of complete metric spaces are complete.

(iv)~A density operator $\rho$ is an extreme point of $\mathfrak{D}(\HC)$
if and only if $\rho^2 = \rho$, which holds precisely when
$\rho = \ket{\psi}\bra{\psi}$ for some unit vector $\psi$.  For the
forward direction: if $\rho$ has rank $\geq 2$, write
$\rho = \lambda \rho_1 + (1-\lambda)\rho_2$ for distinct density operators
supported on orthogonal subspaces.  Converse: $\ket{\psi}\bra{\psi}$ is
idempotent and hence extreme.
\end{proof}

\begin{lemma}[Spectral Decomposition of Density Operators]
\label{lem:3.5}
Every $\rho \in \mathfrak{D}(\HC)$ admits a spectral decomposition
\[
  \rho = \sum_{k} p_k \ket{\psi_k}\bra{\psi_k},
\]
where $\{p_k\}$ are the eigenvalues with $p_k \geq 0$, $\sum_k p_k = 1$,
and $\{\psi_k\}$ is an orthonormal set of eigenvectors.  The sum converges
in trace norm.
\end{lemma}

\begin{proof}
$\rho$ is a compact, self-adjoint, positive, trace-class operator.  By the
spectral theorem for compact self-adjoint operators, $\rho$ has a countable
set of eigenvalues $\{p_k\}$ with $p_k \geq 0$, accumulating only at $0$,
and an orthonormal eigenbasis.  Trace-class convergence follows from
$\sum_k p_k = \Tr(\rho) = 1 < \infty$.
\end{proof}

% ============================================================
\subsection{Interface Map}
\label{sec:ch3:interface}
% ============================================================

We first define the domain algebra.

\begin{definition}[Ontic Operator Algebra]
\label{def:3.15}
Let $\mu$ be a $\sigma$-finite Borel measure on $(\CR, \mathscr{B}(\CR))$.
The \emph{ontic operator algebra} is
\[
  \CR_{\mathrm{op}} := L^2(\CR, \mu),
\]
and $\CB(\CR_{\mathrm{op}})$ denotes the $C^*$-algebra of bounded
operators on $\CR_{\mathrm{op}}$.
\end{definition}

\begin{remark}
\label{rem:3.2}
The choice of $\mu$ reflects the environmental sampling distribution.
Different choices yield unitarily inequivalent representations; this is
analogous to the choice of vacuum in quantum field theory.  The measure
$\mu$ is not an observable of CIIR but a structural parameter.
\end{remark}

\begin{definition}[Complete Positivity]
\label{def:3.16}
A linear map $\Phi : \CB(\CR_{\mathrm{op}}) \to \CB(\HC)$ is
\emph{completely positive} (CP) if for every $n \in \mathbb{N}$, the
ampliation
\[
  \Phi \otimes \id_n : \CB(\CR_{\mathrm{op}}) \otimes M_n(\mathbb{C})
    \to \CB(\HC) \otimes M_n(\mathbb{C})
\]
maps positive elements to positive elements, where $M_n(\mathbb{C})$ is
the algebra of $n \times n$ complex matrices.
\end{definition}

\begin{definition}[Trace Preservation]
\label{def:3.17}
A linear map $\Phi : \mathcal{T}_1(\CR_{\mathrm{op}}) \to
\mathcal{T}_1(\HC)$ is \emph{trace-preserving} (TP) if
$\Tr(\Phi(\sigma)) = \Tr(\sigma)$ for all
$\sigma \in \mathcal{T}_1(\CR_{\mathrm{op}})$.
\end{definition}

\begin{definition}[Push-Forward on the Operator Algebra]
\label{def:3.17b}
Let $\iota : \CM \hookrightarrow \CR$ be the $C^2$ embedding of the
constraint manifold.  For a $\sigma$-finite Borel measure $\mu$ on $\CR$,
the embedding induces a map on $L^2$ spaces via composition:
$\iota^* : L^2(\CR, \mu) \to L^2(\CM, \iota^*\mu)$ defined by
$\iota^* f := f \circ \iota$.  The \emph{push-forward on the operator
algebra} is the $*$-homomorphism
\[
  \iota_* : \CB(\CR_{\mathrm{op}}) \to \CB(\CR_{\mathrm{op}}),
  \qquad
  \iota_*(A) := (\iota^*)^\dagger \circ A \circ \iota^*,
\]
where $(\iota^*)^\dagger$ is the Hilbert adjoint of $\iota^*$ (extended
by zero on $(\mathrm{ran}\,\iota^*)^\perp$).
\end{definition}

\begin{definition}[Interface Map]
\label{def:3.18}
An \emph{interface map} is a linear map
$\Phi : \CB(\CR_{\mathrm{op}}) \to \CB(\HC)$ satisfying:
\begin{enumerate}
  \item[(I1)] $\Phi$ is completely positive (Definition~\ref{def:3.16}).
  \item[(I2)] $\Phi$ is trace-preserving (Definition~\ref{def:3.17})
    on $\mathcal{T}_1(\CR_{\mathrm{op}})$.
  \item[(I3)] \textbf{Constraint covariance:}
    $\Phi \circ \iota_* = \hatP_\CM \circ \Phi$, where $\iota_*$ is the
    push-forward (Definition~\ref{def:3.17b}) and $\hatP_\CM$ is the
    orthogonal projection onto the constraint-image subspace $\HC_\CM$
    (Definition~\ref{def:3.21}).
\end{enumerate}
Condition (I3) is a \emph{consistency requirement}: $\Phi$ is first
constructed satisfying (I1)--(I2), then $\HC_\CM$ is derived from $\Phi$
(Definition~\ref{def:3.21}), and finally (I3) is verified as a constraint
on the admissible class of interface maps.
\end{definition}

\begin{proposition}[Kraus Representation]
\label{prop:3.4}
Every interface map $\Phi$ admits a Kraus decomposition:
\[
  \Phi(\sigma) = \sum_{k=1}^{K} V_k \sigma V_k^\dagger,
  \qquad
  \sum_{k=1}^{K} V_k^\dagger V_k = \id_{\CR_{\mathrm{op}}},
\]
where $V_k : \CR_{\mathrm{op}} \to \HC$ are bounded operators.  If
$\dim\HC = d < \infty$, then $K \leq d^2$.
\end{proposition}

\begin{proof}
By the Stinespring dilation theorem, every CPTP map from a $C^*$-algebra
to $\CB(\HC)$ has an isometric dilation
$\Phi(\sigma) = V \pi(\sigma) V^*$ where $\pi$ is a $*$-representation
on an ancillary Hilbert space $\mathcal{K}$ and $V : \HC \to \mathcal{K}$
is an isometry.  Choosing an orthonormal basis $\{f_k\}$ of the minimal
dilation space and defining $V_k := \langle f_k | V$ yields the Kraus
form.  The Choi--Kraus bound gives $K \leq d^2$ when
$\dim\HC = d < \infty$.
\end{proof}

\begin{lemma}[Interface Maps are Contractive]
\label{lem:3.6}
Every interface map $\Phi$ is a contraction in the trace norm:
$\|\Phi(\sigma)\|_1 \leq \|\sigma\|_1$ for all
$\sigma \in \mathcal{T}_1(\CR_{\mathrm{op}})$.
\end{lemma}

\begin{proof}
For any CPTP map $\Phi$ and any $\sigma \in \mathcal{T}_1$, the
data-processing inequality for the trace norm gives
$\|\Phi(\sigma_1) - \Phi(\sigma_2)\|_1 \leq \|\sigma_1 - \sigma_2\|_1$.
Setting $\sigma_2 = 0$ and using linearity yields the claim.
\end{proof}

\begin{proposition}[Stinespring Dilation of Interface Map]
\label{prop:3.5}
Every interface map $\Phi$ admits a Stinespring dilation: there exist a
Hilbert space $\mathcal{K}$ (the \emph{ancillary environment space}), an
isometry $V : \HC \to \HC \otimes \mathcal{K}$, and a $*$-homomorphism
$\pi : \CB(\CR_{\mathrm{op}}) \to \CB(\HC \otimes \mathcal{K})$ such that
\[
  \Phi(\sigma) = \Tr_{\mathcal{K}}(V \pi(\sigma) V^*).
\]
The environment $\mathcal{K}$ represents the degrees of freedom lost in
the constraint projection.
\end{proposition}

\begin{proof}
This is the Stinespring dilation theorem (Stinespring, 1955).  The
existence of $(\mathcal{K}, V, \pi)$ follows from complete positivity
of $\Phi$.  The partial trace over $\mathcal{K}$ recovers $\Phi$ from
the dilation.
\end{proof}

% ============================================================
\subsection{Observable Algebra}
\label{sec:ch3:observable}
% ============================================================

\begin{definition}[Self-Adjoint Operator]
\label{def:3.19}
An operator $A \in \CB(\HC)$ is \emph{self-adjoint} if $A = A^\dagger$,
where $A^\dagger$ is defined by
$\langle A^\dagger \psi | \phi \rangle = \langle \psi | A\phi \rangle$
for all $\psi, \phi \in \HC$.  The set of bounded self-adjoint operators
is $\CB(\HC)_{\mathrm{sa}} := \{ A \in \CB(\HC) : A = A^\dagger \}$.
\end{definition}

\begin{definition}[Observable Algebra]
\label{def:3.20}
The \emph{observable algebra} of CIIR is the von Neumann algebra
\[
  \CA(\HC) := \bigl\{
    \hatO \in \CB(\HC)_{\mathrm{sa}} :
    [\hatO, \hatP_\CM] = 0
  \bigr\},
\]
where $[\hatO, \hatP_\CM] := \hatO\hatP_\CM - \hatP_\CM\hatO$ is the
commutator and $\hatP_\CM$ is the orthogonal projection onto $\HC_\CM$.
\end{definition}

\begin{proposition}[Algebraic Properties of $\CA(\HC)$]
\label{prop:3.6}
\begin{enumerate}
  \item[(i)] $\CA(\HC)$ is a real vector space closed under the Jordan
    product $A \circ B := \frac{1}{2}(AB + BA)$.
  \item[(ii)] $\CA(\HC) \subseteq \{\hatP_\CM\}'$, the commutant of
    $\hatP_\CM$ in $\CB(\HC)$.
  \item[(iii)] $\CA(\HC)$ is closed in the weak operator topology and
    hence is a von Neumann algebra (as a real $*$-subalgebra of
    $\{\hatP_\CM\}'$).
  \item[(iv)] The identity operator $\id \in \CA(\HC)$.
\end{enumerate}
\end{proposition}

\begin{proof}
(i)~If $A, B \in \CA(\HC)$, then $(A \circ B)^\dagger = A \circ B$
(both are self-adjoint) and
$[A \circ B, \hatP_\CM] = \frac{1}{2}([AB, \hatP_\CM] + [BA, \hatP_\CM])
= \frac{1}{2}(A[B, \hatP_\CM] + [A, \hatP_\CM]B + B[A, \hatP_\CM]
+ [B, \hatP_\CM]A) = 0$ since $[A, \hatP_\CM] = [B, \hatP_\CM] = 0$.

(ii)~By definition, every element of $\CA(\HC)$ commutes with $\hatP_\CM$.

(iii)~The commutant $\{\hatP_\CM\}'$ is a von Neumann algebra by
von~Neumann's double commutant theorem.  The self-adjoint part is closed
in the weak operator topology (the weak limit of self-adjoint operators
is self-adjoint).

(iv)~$\id = \id^\dagger$ and $[\id, \hatP_\CM] = 0$.
\end{proof}

\begin{lemma}[Spectral Theorem for Observables]
\label{lem:3.7}
Every $\hatO \in \CA(\HC)$ admits a spectral decomposition
\[
  \hatO = \int_{\mathbb{R}} \lambda \; dE_\lambda,
\]
where $\{E_\lambda\}_{\lambda \in \mathbb{R}}$ is a projection-valued
measure on $\mathbb{R}$ and $E_\lambda \in \{\hatP_\CM\}'$ for all
$\lambda$.
\end{lemma}

\begin{proof}
The spectral theorem for bounded self-adjoint operators on a Hilbert
space gives the decomposition $\hatO = \int \lambda \; dE_\lambda$.
Since $\hatO \in \{\hatP_\CM\}'$, and the spectral projections
$E_\lambda$ belong to the von Neumann algebra generated by $\hatO$
(which is contained in $\{\hatP_\CM\}'$), we have
$E_\lambda \in \{\hatP_\CM\}'$ for all $\lambda$.
\end{proof}

% ============================================================
\subsection{Constraint-Image Subspace}
\label{sec:ch3:subspace}
% ============================================================

\begin{definition}[Constraint-Image Subspace]
\label{def:3.21}
For an interface map $\Phi$ (Definition~\ref{def:3.18}) and constraint
manifold $\CM$ (Definition~\ref{def:3.7}), the \emph{constraint-image
subspace} is
\[
  \HC_\CM := \overline{\mathrm{span}}\bigl\{
    \Phi(A)\psi : A \in \CB(\CR_{\mathrm{op}}),\; \psi \in \HC
  \bigr\} \subseteq \HC.
\]
The \emph{constraint projection} $\hatP_\CM : \HC \to \HC$ is the
orthogonal projection onto $\HC_\CM$.
\end{definition}

\begin{proposition}[Properties of the Constraint-Image Subspace]
\label{prop:3.7}
\begin{enumerate}
  \item[(i)] $\HC_\CM$ is a closed subspace of $\HC$.
  \item[(ii)] $\hatP_\CM = \hatP_\CM^\dagger = \hatP_\CM^2$
    (self-adjoint idempotent).
  \item[(iii)] $\HC = \HC_\CM \oplus \HC_\CM^\perp$ (orthogonal
    complement decomposition).
  \item[(iv)] For any $\rho \in \mathfrak{D}(\HC)$, the
    \emph{constraint-projected state} $\rho_\CM := \hatP_\CM \rho
    \hatP_\CM / \Tr(\hatP_\CM \rho \hatP_\CM)$ lies in
    $\mathfrak{D}(\HC_\CM)$ whenever $\Tr(\hatP_\CM \rho \hatP_\CM) > 0$.
\end{enumerate}
\end{proposition}

\begin{proof}
(i)~$\HC_\CM$ is defined as the closure of a linear span in a Hilbert
space, which is a closed subspace by definition.

(ii)~The orthogonal projection onto a closed subspace of a Hilbert space
is self-adjoint and idempotent; this is the projection theorem.

(iii)~Every closed subspace of a Hilbert space has an orthogonal
complement, and $\HC = \HC_\CM \oplus \HC_\CM^\perp$ is the orthogonal
decomposition theorem.

(iv)~$\hatP_\CM \rho \hatP_\CM$ is self-adjoint (since $\hatP_\CM$ and
$\rho$ are), positive ($\langle\psi|\hatP_\CM\rho\hatP_\CM|\psi\rangle
= \langle\hatP_\CM\psi|\rho|\hatP_\CM\psi\rangle \geq 0$), and
trace-class (as $\hatP_\CM \rho \hatP_\CM \leq \rho$ in the operator
order).  Normalising by $\Tr(\hatP_\CM\rho\hatP_\CM) > 0$ gives a valid
density operator with support in $\HC_\CM$.
\end{proof}

\begin{lemma}[Dimension Bound on Constraint-Image Subspace]
\label{lem:3.8}
If $\Phi$ has finite Kraus rank $K$ and $\dim\HC < \infty$, then
$\dim\HC_\CM \leq \dim\HC$.  If additionally $\Phi$ is surjective onto
$\CB(\HC_\CM)$, then $\dim\HC_\CM \geq 1$.
\end{lemma}

\begin{proof}
$\HC_\CM \subseteq \HC$ directly gives $\dim\HC_\CM \leq \dim\HC$.
If $\Phi$ is surjective onto $\CB(\HC_\CM)$, then $\HC_\CM$ contains
at least one non-zero vector (the image of any non-zero operator applied
to any non-zero vector), hence $\dim\HC_\CM \geq 1$.
\end{proof}

% ============================================================
\subsection{Consistency Check}
\label{sec:ch3:consistency}
% ============================================================

\noindent\textbf{Verification.}
\begin{enumerate}
  \item \textbf{No undefined symbols:}
    Every symbol in Definitions~\ref{def:3.1}--\ref{def:3.21} is either
    a standard mathematical object ($\mathbb{R}$, $\mathbb{C}$,
    $\sigma$-algebra, etc.)\ or is defined earlier in this chapter.  The
    dependency order is:
    $\CR \to \mathscr{B}(\CR) \to \CR_{\mathrm{op}} \to \CM \to \HC \to
    \CB(\HC) \to \mathcal{T}_1(\HC) \to \mathfrak{D}(\HC) \to \Phi \to
    \CA(\HC) \to \HC_\CM \to \hatP_\CM$.

  \item \textbf{No circular definitions:}
    The dependency graph is a directed acyclic graph.  In particular:
    \begin{itemize}
      \item $\CR$ (Def.~\ref{def:3.1}) depends on no CIIR-specific objects.
      \item $\CM$ (Def.~\ref{def:3.7}) depends on $\CR$ (via embedding $\iota$).
      \item $\HC$ (Def.~\ref{def:3.11}) depends on no prior CIIR objects.
      \item $\Phi$ (Def.~\ref{def:3.18}) is constructed in two phases:
        (a)~conditions (I1)--(I2) depend only on $\CR_{\mathrm{op}}$ and
        $\HC$; (b)~condition (I3) is a consistency requirement verified
        \emph{after} $\HC_\CM$ (Def.~\ref{def:3.21}) is derived from
        $\Phi$.  Thus the construction order is
        $\Phi_{\text{(I1--I2)}} \to \HC_\CM \to \text{verify (I3)}$,
        which is acyclic.
    \end{itemize}

  \item \textbf{No contradiction with prior chapters:}
    Chapter~1 (Introduction) and Chapter~2 (Motivation) contain no formal
    definitions or axioms.  All objects here are introduced for the first
    time.

  \item \textbf{Sufficiency for axiom system:}
    The six primitives ($\CR$, $\CM$, $\HC$, $\Phi$, $\CA(\HC)$,
    $\HC_\CM$) provide all objects referenced in the six CIIR axioms to
    be stated in Chapter~4.
\end{enumerate}

% ============================================================
\subsection{Dependency Graph}
\label{sec:ch3:depgraph}
% ============================================================

The complete dependency structure of Chapter~3 primitives:

\begin{center}
\begin{tabular}{lll}
\hline
\textbf{Object} & \textbf{Definition} & \textbf{Depends on} \\
\hline
$\CR$ & Def.~\ref{def:3.1} & (none) \\
$\mathscr{B}(\CR)$ & Def.~\ref{def:3.2} & $\CR$ \\
$(\CM, g)$ & Defs.~\ref{def:3.3}--\ref{def:3.4} & (none) \\
$\nabla$ & Def.~\ref{def:3.5} & $(\CM, g)$ \\
$d\mu_g$ & Def.~\ref{def:3.6} & $(\CM, g)$ \\
$(\CM, g, \iota, \hatC)$ & Def.~\ref{def:3.7} & $\CR$, $(\CM, g)$ \\
$\kappa$ & Def.~\ref{def:3.8} & $(\CM, g, \iota, \hatC)$, $\nabla$ \\
$F^{\hatC}$ & Def.~\ref{def:3.9} & $\hatC$, $\nabla$ \\
$\HC$ & Defs.~\ref{def:3.10}--\ref{def:3.11} & (none) \\
$\CB(\HC)$ & Def.~\ref{def:3.12} & $\HC$ \\
$\mathcal{T}_1(\HC)$ & Def.~\ref{def:3.13} & $\HC$, $\CB(\HC)$ \\
$\mathfrak{D}(\HC)$ & Def.~\ref{def:3.14} & $\mathcal{T}_1(\HC)$ \\
$\CR_{\mathrm{op}}$ & Def.~\ref{def:3.15} & $\CR$, $\mathscr{B}(\CR)$ \\
$\iota_*$ & Def.~\ref{def:3.17b} & $\iota$, $\CR_{\mathrm{op}}$ \\
$\Phi$ (construction) & Def.~\ref{def:3.18}(I1--I2) & $\CR_{\mathrm{op}}$, $\HC$ \\
$\Phi$ (covariance) & Def.~\ref{def:3.18}(I3) & $\iota_*$, $\hatP_\CM$ (verified post-construction) \\
$\CA(\HC)$ & Def.~\ref{def:3.20} & $\CB(\HC)$, $\hatP_\CM$ \\
$\HC_\CM$ & Def.~\ref{def:3.21} & $\Phi$, $\HC$ \\
$\hatP_\CM$ & Def.~\ref{def:3.21} & $\HC_\CM$ \\
\hline
\end{tabular}
\end{center}

\noindent
All 21 definitions, 8 lemmas/propositions, and their proofs are complete.
The primitive vocabulary is sufficient to state the six CIIR axioms in
Chapter~4.
